\chapter{Weil--Petersson Kernel Comparison: Dense Test-Class Limits}
\label{v4-arith:w6-b:chapter}

\emph{The bounded-interval reduction, construction B. Chapter~\ref{v4-arith:w5-a:chapter} identifies
the kernel mismatch: Rankin--Selberg Petersson norms give positive
automorphic kernels, but the zeta finite-prime Weil term is the
von Mangoldt distribution. Four natural
extensions of the automorphic kernels widen the
comparison class and clarify the topology, but they do not prove
\(W_{\mathrm{fin}}(f*\widetilde f)\ge0\). The full residual remains the
A-TP inequality of
\Cref{v4-arith:w5-a:conj:archimedean-tate-positivity}, with finite,
polar, and archimedean terms all present.}

\section{Restricted-Class Baseline and Residual}
\label{v4-arith:w6-b:baseline}

Fix notation. \(\mathrm{PW}(\mathbb{R})\) is the
Paley--Wiener class of
\Cref{v4-arith:w5-a:def:PW}: even Schwartz-class functions of compact
support on \(\mathbb{R}\) with Fourier transforms that are entire of
exponential type. For \(f\in\mathrm{PW}(\mathbb{R})\), write
\(\widetilde{f}(u)=\overline{f(-u)}\), so \(f\) real even forces
\(\widetilde{f}=f\). Write \(W(\cdot)\) for the Weil distribution
\eqref{v4-arith:w5-a:eq:weil-explicit}, decomposed into polar terms,
\(W_{\infty}\) (archimedean Tate), and \(W_{\mathrm{fin}}\) (finite-prime).

\begin{theorem}[the restricted-class comparison baseline kernel comparison]
\label{v4-arith:w6-b:thm:baseline}
\ClaimStatusProvedHere. For a normalised holomorphic cuspidal Hecke
eigenform \(F\), the coefficientwise lift \(\Upsilon_F\) is a formal
automorphic probe. Its Rankin--Selberg kernel is governed by
\(L(s,F\otimes\overline F)\). It is not the finite-prime Weil
distribution \(W_{\mathrm{fin}}\) of \(\zeta\).
\end{theorem}

\begin{proof}
\Cref{v4-arith:w5-a:thm:hpli-form} and
\Cref{v4-arith:w5-a:thm:prime-side}.
\end{proof}

\begin{remark}[domain boundary named precisely]
\label{v4-arith:w6-b:rem:domain-boundary}
Three obstructions in \Cref{v4-arith:w5-a:rem:H2-obstruction} and
\Cref{v4-arith:w5-a:prop:eisenstein-residual} structure the gap
between \(\mathrm{PW}^{\mathrm{hol}}_{F}\) and the full \(\mathrm{PW}(\mathbb{R})\).
\begin{itemize}
\item[\textbf{(O1)}] \emph{Mellin-strip width.} Hypothesis (H2)
restricts \(\widetilde{f}_{M}(s)=\int_{0}^{\infty}f(\log u)u^{s-1}du\)
to the strip \(|\mathrm{Re}(s)-1/2|\leq (k-1)/2-\varepsilon\). A fixed
\(F\) of finite weight \(k\) cannot accommodate \(f\) whose Mellin strip is
wider.
\item[\textbf{(O2)}] \emph{Holomorphic archimedean type.} The lift
\(\Upsilon_{F}\) uses the holomorphic discrete series at \(\infty\) and
sees only holomorphic cuspidal representations. Maass cusp forms of
Laplace eigenvalue \(\lambda=1/4+r^{2}\) are not in the image.
\item[\textbf{(O3)}] \emph{Discrete-spectrum restriction.} The
isometry lands in the cuspidal core, orthogonal to the residual and
the Eisenstein continuous spectrum. Part of the Weil distribution
sees these other sectors through the Mellin transform, even though
the non-trivial zero sum is discrete-only
(\Cref{v4-arith:w5-c:prop:domain-suffices-for-Li}).
\end{itemize}
Obstructions (O1)--(O3) are three orthogonal directions in which the
restricted class \(\mathrm{PW}^{\mathrm{hol}}_{F}\) fails to exhaust
\(\mathrm{PW}(\mathbb{R})\); each admits a distinct extension construction,
and the four extension constructions below address them jointly.
\end{remark}

\section{Construction I: Maass-Form Lift on Tempered Eigenforms}
\label{v4-arith:w6-b:construction-I}

\subsection{Archimedean Local Factor for Principal-Series Maass}
\label{v4-arith:w6-b:construction-I:local}

Let \(\phi\) be a Maass cuspidal Hecke eigenform on
\(\mathrm{SL}_{2}(\mathbb{Z})\backslash\mathbb{H}\) with Laplace
eigenvalue \(\lambda_{\phi}=1/4+r_{\phi}^{2}\), \(r_{\phi}\in\mathbb{R}\)
tempered at \(\infty\). The Fourier--Whittaker expansion is
\begin{equation}
\label{v4-arith:w6-b:eq:maass-expansion}
\phi(z)
\;=\;\sum_{n\neq 0}a_{\phi}(n)\,\sqrt{y}\,K_{ir_{\phi}}(2\pi|n|y)\,e^{2\pi inx},
\end{equation}
where \(K_{\nu}\) is the Macdonald \(K\)-Bessel function (Iwaniec 2002
\emph{Spectral Methods of Automorphic Forms} \S1.9). The normalised
Hecke eigenvalues \(\widetilde{a}_{\phi}(n)=a_{\phi}(n)/\sqrt{n}\) satisfy
\(\widetilde{a}_{\phi}(1)=1\) and are real if \(\phi\) is a newform at
level~\(1\). Deligne's bound is not available for Maass forms; the
current best is Kim--Sarnak \(\theta_{\phi}\leq 7/64\), and temperedness
at every prime is the Ramanujan conjecture, open at
\(\mathrm{GL}_{2}/\mathbb{Q}\).

\begin{definition}[Maass Mellin lift]
\label{v4-arith:w6-b:def:upsilon-maass}
\ClaimStatusDefinitional. Fix a tempered Maass Hecke eigenform \(\phi\)
of eigenvalue \(\lambda=1/4+r^{2}\) with \(r\in\mathbb{R}\). For
\(f\in\mathrm{PW}(\mathbb{R})\), put
\begin{equation}
\label{v4-arith:w6-b:eq:upsilon-maass}
\Upsilon_{\phi}(f)(z)
\;:=\;\sum_{n\geq 1}f(\log n)\,\widetilde{a}_{\phi}(n)\,\sqrt{y}\,K_{ir}(2\pi ny)\cos(2\pi nx).
\end{equation}
The cosine is the symmetrisation under \(\phi\) even (the \(+\) part of
the Maass space; the odd case replaces cosine by sine and the
archimedean Whittaker by its derivative).
\end{definition}

\begin{lemma}[\(\Upsilon_{\phi}(f)\) is a Maass-type function on the tempered class]
\label{v4-arith:w6-b:lem:upsilon-maass-cusp}
\ClaimStatusProvedHere. Let
\begin{equation}
\label{v4-arith:w6-b:eq:PW-maass}
\mathrm{PW}^{\mathrm{Maass}}_{\phi}\;\subset\;\mathrm{PW}(\mathbb{R})
\end{equation}
denote the subspace of \(f\in\mathrm{PW}(\mathbb{R})\) satisfying:
\begin{itemize}
\item[(M1)] \(f\) is real even;
\item[(M2)] the Mellin transform \(\widetilde{f}_{M}\) is holomorphic
and of Schwartz decay on the strip
\(|\mathrm{Re}(s)-1/2|\leq 1/2-7/64-\varepsilon\) for some
\(\varepsilon>0\) (Kim--Sarnak width);
\item[(M3)] \(\widetilde{f}_{M}\) has a double zero at \(s=0\) and \(s=1\).
\end{itemize}
If \(\phi\) is tempered at every finite prime (\(\theta_{p}=0\)), then
\(\Upsilon_{\phi}(f)\) is a Maass-type cuspidal function on
\(\mathrm{SL}_{2}(\mathbb{Z})\backslash\mathbb{H}\) with Laplace
eigenvalue \(\lambda_{\phi}\) on the restricted spectral locus, and its
Fourier--Whittaker coefficients are
\(f(\log n)\widetilde{a}_{\phi}(n)\), absolutely summable on
\(\mathrm{PW}^{\mathrm{Maass}}_{\phi}\) under the Kim--Sarnak bound.
\end{lemma}

\begin{proof}
Hypothesis (M2) gives the convergence strip; (M3) removes the Tate
polar contribution. Under \(\theta_{p}=0\),
\(|\widetilde{a}_{\phi}(n)|\leq d(n)\) by the Ramanujan hypothesis
applied place-by-place, and the series converges absolutely by
Schwartz decay of \(\widetilde{f}_{M}\) combined with subconvex Hecke
eigenvalue growth. Modularity and cuspidality follow from the
functional equation of \(L(s,\phi)\) (Jacquet--Langlands SLN~114 Thm.~11.1)
combined with the Maass--Selberg formalism
(Moeglin--Waldspurger 1995 II.1).
Unconditionally (without Ramanujan for \(\phi\)), replace (M2) by
\(|\mathrm{Re}(s)-1/2|\leq 1/2-7/64-\varepsilon\) from Kim--Sarnak; the
convergence and cuspidality survive, at the cost of a narrower strip.
\end{proof}

\begin{theorem}[Maass Rankin--Selberg isometry]
\label{v4-arith:w6-b:thm:maass-RS}
\ClaimStatusProvedHere \emph{(domain-restricted to
\(\mathrm{PW}^{\mathrm{Maass}}_{\phi}\) for \(\phi\) tempered at every
prime)}. For \(f\in\mathrm{PW}^{\mathrm{Maass}}_{\phi}\), the Petersson
norm of \(\Upsilon_{\phi}(f)\) on
\(\mathrm{SL}_{2}(\mathbb{Z})\backslash\mathbb{H}\) equals the
Rankin--Selberg integral
\begin{equation}
\label{v4-arith:w6-b:eq:maass-RS}
\bigl\langle\Upsilon_{\phi}(f),\Upsilon_{\phi}(f)\bigr\rangle^{\mathrm{Pet}}
\;=\;\int_{-\infty}^{\infty}|\widehat{f}(t)|^{2}\cdot K_{\phi}\bigl(\tfrac{1}{2}+it\bigr)\,dt,
\end{equation}
where
\begin{equation}
\label{v4-arith:w6-b:eq:Kphi}
K_{\phi}(s)\;:=\;\frac{\zeta(2s)}{\zeta(s)\,\xi_{\mathbb{R}}(s)\xi_{\mathbb{R}}(1-s)}\cdot L(s,\phi\otimes\overline{\phi})\cdot G_{\phi}(s),
\end{equation}
with \(L(s,\phi\otimes\overline{\phi})\) the Rankin--Selberg
\(L\)-function of Jacquet--Piatetski-Shapiro--Shalika (JPSS~1983) and
\(G_{\phi}(s)=|\Gamma(\tfrac{s}{2}+ir)\Gamma(\tfrac{s}{2}-ir)|/|\Gamma(\tfrac{s}{4}+\tfrac{ir}{2})\Gamma(\tfrac{s}{4}-\tfrac{ir}{2})|^{2}\)
the archimedean Whittaker absorption factor from the \(K\)-Bessel
identity
\(\int_{0}^{\infty}K_{ir}(y)^{2}y^{s-1}dy=G_{\phi}(s)\)
(Gradshteyn--Ryzhik 6.576.4). The kernel \(K_{\phi}\) is non-negative on
the critical line \(\mathrm{Re}(s)=1/2\) on the tempered locus.
\end{theorem}

\begin{proof}
Unfolding the Petersson integral of \(|\Upsilon_{\phi}(f)|^{2}\) on
\(\mathrm{SL}_{2}(\mathbb{Z})\backslash\mathbb{H}\) against the
fundamental domain gives
\begin{equation*}
\langle\Upsilon_{\phi}(f),\Upsilon_{\phi}(f)\rangle^{\mathrm{Pet}}
=\sum_{m,n}f(\log m)\overline{f(\log n)}\widetilde{a}_{\phi}(m)\overline{\widetilde{a}_{\phi}(n)}\int\! K_{ir}(2\pi my)K_{ir}(2\pi ny)y\,dy\cdot\delta_{m,n}.
\end{equation*}
The \(x\)-integral of \(\cos(2\pi mx)\cos(2\pi nx)\) over
\([-1/2,1/2]\) gives \(\delta_{m,n}/2\). The \(y\)-integral is the
regularised Rankin--Selberg local zeta integral, not the divergent bare
integral of \(K_{ir}(y)^{2}dy/y\); in the normalisation of Zagier 1981
\S I.3 and Iwaniec 2002 Thm.~10.2 it contributes a positive archimedean
Plancherel factor. The diagonal sum becomes
\begin{equation*}
=C_{\phi}\sum_{n\geq 1}\frac{|f(\log n)|^{2}|\widetilde{a}_{\phi}(n)|^{2}}{n}
=C_{\phi}\cdot\frac{1}{(2\pi i)^{2}}\int\!\int|\widetilde{f}_{M}(s)|^{2}L(s,\phi\otimes\overline{\phi})/\zeta(2s)\,ds\,d\overline{s},
\end{equation*}
where \(C_{\phi}\) is the archimedean absorption constant and the
inner Dirichlet series identity is the Rankin--Selberg formula for
Maass forms (Rankin 1939 for holomorphic; Bump 1997 \S3.7 for Maass).
Identification with \eqref{v4-arith:w6-b:eq:maass-RS} uses
\(\widetilde{f}_{M}(1/2+it)=\widehat{f}(t)\) up to normalisation and
absorbs \(\xi_{\mathbb{R}}\) and \(G_{\phi}\) into the archimedean
factor of \(K_{\phi}\). Non-negativity: the Rankin--Selberg unfolding
expresses the Petersson norm as
\(C_{\phi}\sum_{n\geq 1}|f(\log n)|^{2}|\widetilde{a}_{\phi}(n)|^{2}/n\),
a sum of non-negative terms; hence \(K_{\phi}(1/2+it)\geq 0\) for all
\(t\in\mathbb{R}\) on the tempered locus. Ramanujan controls the
tempered local bounds; the positivity comes from the squared
Rankin--Selberg norm.
\end{proof}

\begin{corollary}[Maass construction gives automorphic-kernel positivity, not zeta prime-side closure]
\label{v4-arith:w6-b:cor:maass-closure}
\ClaimStatusProvedHereConditional. Under Ramanujan for a fixed Maass eigenform
\(\phi\), the Rankin--Selberg norm attached to \(\phi\) is positive on
its automorphic kernel. This does not prove
\(W_{\mathrm{fin}}(f*\widetilde f)\ge0\), because the zeta finite-place
term has von Mangoldt coefficients, not the diagonal coefficients of
\(\phi\).
\end{corollary}

\begin{proof}
Rankin--Selberg unfolding gives a positive Petersson norm with
coefficients \(|a_{\phi}(n)|^2\). Weil's zeta finite term is the
logarithmic derivative of \(\zeta\), with coefficients \(\Lambda(n)\).
These distributions are not identified by Ramanujan or by the
Kim--Sarnak bound.
\end{proof}

\begin{remark}[union over Maass eigenforms]
\label{v4-arith:w6-b:rem:maass-union}
The full Maass cuspidal spectrum of
\(\mathrm{SL}_{2}(\mathbb{Z})\backslash\mathbb{H}\) is
\(\{\phi_{j}\}_{j\geq 1}\) with eigenvalues
\(\lambda_{j}=1/4+r_{j}^{2}\) accumulating at infinity with Weyl-law
density \(\#\{j:|r_{j}|\leq R\}\sim R^{2}/12\) (Selberg 1956). The
per-\(\phi_{j}\) restricted class
\(\mathrm{PW}^{\mathrm{Maass}}_{\phi_{j}}\) depends on \(\phi_{j}\) only
through the temperedness hypothesis (and the archimedean factor
\(G_{\phi_{j}}\), which modulates the spectral density on the critical
line but not the holomorphy strip). The union
\(\bigcup_{j}\mathrm{PW}^{\mathrm{Maass}}_{\phi_{j}}\) collapses to
\(\mathrm{PW}^{\mathrm{Maass}}_{*}\), the common intersection strip of
width \(1/2-7/64-\varepsilon\) unconditionally (or width \(1/2\) under
Ramanujan for all Maass forms), whose union with
\(\mathrm{PW}^{\mathrm{hol}}_{F}\) still falls short of the full
\(\mathrm{PW}(\mathbb{R})\); (O1) is narrowed from width \((k-1)\)
to width slightly under \(1/2\), which is a strict widening at
\(k=1\) but a strict \emph{narrowing} at \(k\geq 2\).
\end{remark}

\section{Construction II: Weight \(k\to\infty\) Limit Along Compatible Newforms}
\label{v4-arith:w6-b:construction-II}

\subsection{Compatible Newform Families}
\label{v4-arith:w6-b:construction-II:family}

The Mellin-strip width \((k-1)\) of the holomorphic construction grows without
bound as the weight \(k\to\infty\); a compatible family of newforms
of increasing weight furnishes an exhausting sequence of restricted
classes whose union is dense in \(\mathrm{PW}(\mathbb{R})\).

\begin{definition}[compatible family of newforms]
\label{v4-arith:w6-b:def:compatible-family}
\ClaimStatusDefinitional. A \emph{compatible family} of holomorphic
cuspidal newforms is a sequence \(\{F_{k}\}_{k\geq k_{0}}\) with
\(F_{k}\in S_{k}(\mathrm{SL}_{2}(\mathbb{Z}))\) newform of weight \(k\),
normalised Hecke eigenvalues \(\widetilde{a}_{F_{k}}(n)\) with
\(\widetilde{a}_{F_{k}}(1)=1\), and no weight-restriction other than
\(k_{0}\geq 12\).
\end{definition}

\begin{lemma}[strip width grows with \(k\)]
\label{v4-arith:w6-b:lem:strip-width}
\ClaimStatusProvedHere. For each \(k\geq 12\), the restricted class
\(\mathrm{PW}^{\mathrm{hol}}_{F_{k}}\) of \Cref{v4-arith:w5-a:lem:upsilon-cuspidal}
admits Mellin-strip width \((k-1)\). As \(k\to\infty\), the strip fills
the whole plane: \(\bigcup_{k}\{\widetilde{f}_{M}\text{ holomorphic on
the strip of width \((k-1)\)}\}=\{\widetilde{f}_{M}\text{ entire}\}\).
\end{lemma}

\begin{proof}
Immediate from (H2): the strip is exactly
\(|\mathrm{Re}(s)-1/2|\leq(k-1)/2-\varepsilon\), which widens
monotonically in \(k\) and fills \(\mathbb{C}\) in the limit.
\end{proof}

\begin{proposition}[union density in Paley--Wiener, subdense weight-wise]
\label{v4-arith:w6-b:prop:union-density-weight}
\ClaimStatusProvedHere. The union
\begin{equation}
\label{v4-arith:w6-b:eq:union-weight}
\mathrm{PW}^{\mathrm{hol}}_{\infty}\;:=\;\bigcup_{k\geq 12}\mathrm{PW}^{\mathrm{hol}}_{F_{k}}
\end{equation}
over a compatible family of newforms is dense in \(\mathrm{PW}(\mathbb{R})\)
in the Schwartz topology of Paley--Wiener functions: for every
\(f\in\mathrm{PW}(\mathbb{R})\) there exists a sequence
\(\{f_{k}\}\subset\mathrm{PW}^{\mathrm{hol}}_{\infty}\) with
\(f_{k}\to f\) in all Schwartz seminorms.
\end{proposition}

\begin{proof}
Given \(f\in\mathrm{PW}(\mathbb{R})\) with \(\widetilde{f}_{M}(s)\)
entire of some finite order (the Paley--Wiener class imposes
exponential type in \(s\)-Mellin dual), choose \(k\) large enough that
the strip of width \((k-1)\) contains all essential Mellin mass of
\(\widetilde{f}_{M}\) up to truncation error \(<1/k\). Define
\(f_{k}\in\mathrm{PW}^{\mathrm{hol}}_{F_{k}}\) by truncating
\(\widetilde{f}_{M}\) to the strip and inverting the Mellin transform;
the truncation error vanishes in all Schwartz seminorms as
\(k\to\infty\) by the Paley--Wiener regularity of \(f\). Hypothesis (H3)
(Tate polar vanishing) is handled by subtracting a finite-rank Taylor
correction of \(\widetilde{f}_{M}\) at \(s=0,1\)
(\Cref{v4-arith:w5-a:rem:H3-removable}); this correction is uniformly
bounded in \(k\) and preserves the Schwartz topology.
\end{proof}

\subsection{The Petersson Measure Degenerates at \(k\to\infty\)}
\label{v4-arith:w6-b:construction-II:petersson-limit}

The weight limit is not cost-free. The Petersson measure on
\(S_{k}(\mathrm{SL}_{2}(\mathbb{Z}))\) degenerates as \(k\to\infty\), and
the Rankin--Selberg kernel \(K_{F_{k}}(s)\) of
\eqref{v4-arith:w5-a:eq:KF-kernel} must be controlled through the
limit.

\begin{proposition}[Rankin--Selberg kernel and the symmetric-square factor]
\label{v4-arith:w6-b:prop:Kk-asymptotic}
\ClaimStatusProvedHere. For a newform \(F_{k}\) of weight \(k\),
\begin{equation}
\label{v4-arith:w6-b:eq:Kk-asymptotic}
K_{F_{k}}\bigl(\tfrac{1}{2}+it\bigr)
\;=\;c_{k}\cdot\frac{\zeta(1+2it)\,\zeta(1-2it)}{\xi_{\mathbb{R}}(\tfrac{1}{2}+it)\xi_{\mathbb{R}}(\tfrac{1}{2}-it)}
\cdot A_{F_{k}}(t),
\end{equation}
where \(c_{k}=\|F_{k}\|^{2}_{\mathrm{Pet}}/[\mathrm{SL}_{2}(\mathbb{Z}):1]\cdot(4\pi)^{-(k-1)}(k-2)!\)
is the Petersson normalising scalar and \(A_{F_{k}}(t)\) is the
normalised symmetric-square Rankin factor. It is built from the
archimedean gamma quotient and \(L(1/2+it,\mathrm{sym}^{2}F_{k})\), with
the usual regularised value at \(t=0\). No theorem quoted here sets
\(A_{F_{k}}(t)=1+O_{t}(k^{-1})\) for an individual form.
\end{proposition}

\begin{proof}
The Rankin--Selberg \(L\)-function \(L(s,F_{k}\otimes\overline{F_{k}})\)
decomposes as \(L(s,F_{k}\otimes\overline{F_{k}})=\zeta(s)L(s,\mathrm{sym}^{2}F_{k})\)
via the \(GL_{2}\)-Clebsch decomposition (Shimura 1975
\emph{Math.~Ann.}~217 \S3). The remaining local calculation is the
standard Rankin--Selberg archimedean factor. After the Petersson scalar
is separated, all dependence on the chosen newform which is not already
visible in the zeta factor is exactly \(A_{F_{k}}(t)\).
\end{proof}

\begin{conjecture}[symmetric-square weight-limit datum]
\label{v4-arith:w6-b:conj:sym2-weight-limit}
\ClaimStatusConjectured. For the chosen compatible family
\(\{F_{k}\}\), the factors \(A_{F_{k}}(t)\) converge to \(1\) locally
uniformly on the real line, with a \(k\)-uniform Schwartz majorant after
multiplication by every fixed Paley--Wiener test function.
\end{conjecture}

\begin{theorem}[weight-limit extension, with Petersson normalisation]
\label{v4-arith:w6-b:thm:weight-limit}
\ClaimStatusProvedHereConditional, conditional on
\Cref{v4-arith:w6-b:conj:sym2-weight-limit}. For
\(f\in\mathrm{PW}(\mathbb{R})\) and
\(\{f_{k}\}\subset\mathrm{PW}^{\mathrm{hol}}_{\infty}\) a Schwartz
approximating sequence with
\(f_{k}\in\mathrm{PW}^{\mathrm{hol}}_{F_{k}}\) (existing by
\Cref{v4-arith:w6-b:prop:union-density-weight}), the
\emph{normalised} Rankin--Selberg norm converges:
\begin{equation}
\label{v4-arith:w6-b:eq:weight-limit}
\lim_{k\to\infty}\frac{1}{c_{k}}\bigl\langle\Upsilon_{F_{k}}(f_{k}),\Upsilon_{F_{k}}(f_{k})\bigr\rangle^{\mathrm{Pet}}
\;=\;\int_{-\infty}^{\infty}|\widehat{f}(t)|^{2}\cdot\frac{\zeta(1+2it)\zeta(1-2it)}{\xi_{\mathbb{R}}(\tfrac{1}{2}+it)\xi_{\mathbb{R}}(\tfrac{1}{2}-it)}\,dt.
\end{equation}
\end{theorem}

\begin{proof}
Combine the Rankin--Selberg obstruction
\Cref{v4-arith:w5-a:thm:RS-integral} for each \(F_{k}\) with the kernel
factorisation \Cref{v4-arith:w6-b:prop:Kk-asymptotic} and the
convergence datum
\Cref{v4-arith:w6-b:conj:sym2-weight-limit}. The Schwartz convergence
\(f_{k}\to f\) of
\Cref{v4-arith:w6-b:prop:union-density-weight} and dominated
convergence on the critical line deliver the limit. The normalising
scalar \(c_{k}\) absorbs the weight-dependence on both sides.
\end{proof}

\begin{remark}[weight-limit is effectively the Selberg trace on \(GL_{2}\)]
\label{v4-arith:w6-b:rem:weight-limit-trace}
Assuming \Cref{v4-arith:w6-b:conj:sym2-weight-limit}, the right-hand
side of \eqref{v4-arith:w6-b:eq:weight-limit} is the
\emph{symmetric Plancherel kernel} at the tempered unitary principal
series of \(\mathrm{GL}_{2}(\mathbb{R})\) integrated against
\(|\widehat{f}|^{2}\). By the \(\mathrm{GL}_{2}(\mathbb{Q})\) Selberg
trace formula (Gelbart 1975 \emph{Automorphic Forms on Adele Groups}
\S6), this kernel is the trace of the convolution
\(f*\widetilde{f}\) against the tempered unitary dual of
\(\mathrm{GL}_{2}(\mathbb{A})\) restricted to holomorphic discrete
series in the archimedean factor. The weight-limit reconstructs the
full holomorphic-discrete-series contribution to the Selberg trace
from its per-weight pieces. Obstruction (O2) (Maass exclusion) is
not lifted by this construction.
\end{remark}

\begin{corollary}[weight-limit preserves automorphic positivity, not \(W_{\mathrm{fin}}\)]
\label{v4-arith:w6-b:cor:weight-limit-positive}
\ClaimStatusProvedHereConditional, conditional on
\Cref{v4-arith:w6-b:conj:sym2-weight-limit}. Under the
Schwartz-limit normalisation of
\Cref{v4-arith:w6-b:thm:weight-limit}, Petersson positivity transfers
to the limiting automorphic Rankin--Selberg kernel. This gives no
positivity statement for \(W_{\mathrm{fin}}\) unless that limiting
kernel is separately identified with the von Mangoldt distribution.
\end{corollary}

\begin{proof}
Weak limits preserve positivity of the automorphic kernels. They do
not change their coefficient system into \(\Lambda(n)\).
\end{proof}

\section{Construction III: Residual-Spectrum Lift and One-Dimensional Absorption}
\label{v4-arith:w6-b:construction-III}

\subsection{The Residual Spectrum at \(GL_{2}/\mathbb{Q}\)}
\label{v4-arith:w6-b:construction-III:residual}

The residual spectrum \(\mathcal{V}^{\mathrm{prim}}_{\mathrm{res}}\) at
\(\mathrm{GL}_{2}/\mathbb{Q}\) is one-dimensional by
Moeglin--Waldspurger 1995 II.2, spanned by the Fock vacuum
\(|0\rangle\), which corresponds to the trivial character twist of the
determinant (\Cref{v4-arith:w5-c:rem:CAP-residual}).

\begin{lemma}[residual Mellin absorption]
\label{v4-arith:w6-b:lem:residual-absorption}
\ClaimStatusProvedHere. The Mellin lift of any
\(f\in\mathrm{PW}(\mathbb{R})\) satisfying the constant-mode
hypothesis
\begin{equation}
\label{v4-arith:w6-b:eq:constant-mode}
\widetilde{f}_{M}(1)=\int_{-\infty}^{\infty}f(u)e^{u}\,du\;=\;0
\end{equation}
is orthogonal to \(|0\rangle\in\mathcal{V}^{\mathrm{prim}}_{\mathrm{res}}\).
For \(f\in\mathrm{PW}(\mathbb{R})\) \emph{without} the constant-mode
hypothesis, the Mellin lift picks up a residual contribution
\(\Upsilon_{\mathrm{res}}(f)=\widetilde{f}_{M}(1)\cdot|0\rangle\) absorbing the
\(t=0\) mode, which is the coefficient of the simple pole of
\(\Lambda_{\zeta}\) at \(s=1\) (Tate 1950).
\end{lemma}

\begin{proof}
The vacuum \(|0\rangle\) pairs with any Mellin transform by
\(\widetilde{f}_{M}(1)=\int_{0}^{\infty}f(\log u)\,du=\int_{-\infty}^{\infty}f(u)e^{u}du\).
This is the Tate-normalised zero mode, not the unweighted additive
Fourier coefficient. The
constant-mode hypothesis kills this pairing; absent the hypothesis,
the coefficient is explicit and absorbs into the polar term
\(\widehat{f*\widetilde{f}}(\pm 1/(2i))\) of the Weil formula.
\end{proof}

\begin{proposition}[residual construction covers the Tate polar]
\label{v4-arith:w6-b:prop:residual-covers-polar}
\ClaimStatusProvedHere. The residual lift
\(\Upsilon_{\mathrm{res}}\) of \Cref{v4-arith:w6-b:lem:residual-absorption},
combined with Petersson positivity at the vacuum
\(\langle|0\rangle,|0\rangle\rangle^{\mathrm{Pet}}=1\), absorbs the polar terms
\(\widehat{f*\widetilde{f}}(\pm 1/(2i))\) of the Weil formula
\eqref{v4-arith:w5-a:eq:weil-explicit} into the isometry side; no
restriction on \(f\) is required beyond \(f\in\mathrm{PW}(\mathbb{R})\).
\end{proposition}

\begin{proof}
Residual + cuspidal decomposes as
\(\mathcal{V}^{\mathrm{prim}}_{\mathrm{disc}}=\mathcal{V}^{\mathrm{prim}}_{\mathrm{cusp}}\oplus\mathcal{V}^{\mathrm{prim}}_{\mathrm{res}}\)
orthogonally (Langlands 1966 \S6); the residual component is
one-dimensional, and its Petersson norm equals \(1\) by Weil's Tamagawa
normalisation (Weil 1967 Ch.~VII; \Cref{v4-arith:kp:lem:adelic-convergence}).
Paired with \Cref{v4-arith:w6-b:lem:residual-absorption}, the polar
coefficient \(\widetilde{f}_{M}(1)\) produces
\(|\widetilde{f}_{M}(1)|^{2}\cdot 1\) on the isometry side, matching the
Weil polar contribution exactly after the \(\xi_{\mathbb{R}}\) polar
normalisation. This replaces Hypothesis (H3)
(\Cref{v4-arith:w5-a:rem:H3-removable}) by the residual sector
absorption, removing Tate-polar vanishing as an input.
\end{proof}

\begin{remark}[Construction III lifts (H3)]
\label{v4-arith:w6-b:rem:construction-III}
Construction~III absorbs the Tate-polar contribution into the
one-dimensional residual sector, eliminating Hypothesis (H3). It
contributes nothing to (O1) or (O2): the residual sector is one
vector, not dense in any sector larger than its own span. The value
of Construction~III is in removing Hypothesis (H3) as a restriction on the
test class.
\end{remark}

\section{Construction IV: Jacquet--Piatetski-Shapiro--Shalika on \(GL_{n}\times GL_{m}\)}
\label{v4-arith:w6-b:construction-IV}

\subsection{Higher-Rank Rankin--Selberg}
\label{v4-arith:w6-b:construction-IV:higher-rank}

The Jacquet--Piatetski-Shapiro--Shalika (JPSS 1983) Rankin--Selberg
integrals for \(GL_{n}\times GL_{m}\) furnish a Mellin-strip width
that grows with the rank \(n\), admitting test classes beyond those
available from \(GL_{2}\).

\begin{definition}[higher-rank Mellin lift, \(GL_{n}\) cusp]
\label{v4-arith:w6-b:def:upsilon-GLn}
\ClaimStatusDefinitional. Fix a rank \(r\geq 2\) and a cuspidal automorphic
representation \(\pi\) of \(\mathrm{GL}_{r}/\mathbb{Q}\). For
\(f\in\mathrm{PW}(\mathbb{R})\), define
\begin{equation}
\label{v4-arith:w6-b:eq:upsilon-GLn}
\Upsilon_{\pi}(f)\;:=\;\sum_{m\geq 1}f(\log m)\,a_{\pi}(m)\,W_{\pi}(m),
\end{equation}
where \(a_{\pi}(m)\) is the Hecke eigenvalue of \(\pi\) at \(m\) (in the
standard Satake normalisation for \(\mathrm{GL}_{r}\)) and \(W_{\pi}(m)\)
is the Whittaker value at \(m\) (via the \(GL_{r}\) Whittaker model,
Shalika 1974).
\end{definition}

\begin{theorem}[higher-rank Rankin--Selberg isometry]
\label{v4-arith:w6-b:thm:higher-rank-RS}
\ClaimStatusProvedHere \emph{(domain-restricted to the
JPSS-convergence cone for \(\pi\))}. For \(\pi\) tempered cuspidal on
\(\mathrm{GL}_{n}/\mathbb{Q}\) and \(f\) in the restricted class
\(\mathrm{PW}^{\mathrm{JPSS}}_{\pi}\) whose Mellin strip sits inside
the JPSS convergence cone for \(\pi\otimes\overline{\pi}\),
\begin{equation}
\label{v4-arith:w6-b:eq:higher-rank-RS}
\bigl\langle\Upsilon_{\pi}(f),\Upsilon_{\pi}(f)\bigr\rangle^{\mathrm{Pet}}
\;=\;\int_{-\infty}^{\infty}|\widehat{f}(t)|^{2}\cdot K_{\pi}\bigl(\tfrac{1}{2}+it\bigr)\,dt,
\end{equation}
with \(K_{\pi}(s)=L(s,\pi\otimes\overline{\pi})/[\zeta(ns)\cdot\prod_{v}G_{v,\pi}(s)]\)
the JPSS kernel (JPSS~1983 \S2). The kernel is non-negative on the
critical line under \(\pi\) tempered everywhere (local unitary
factors).
\end{theorem}

\begin{proof}
JPSS 1983 Thm.~2.7 gives the Rankin--Selberg integral for
\(GL_{n}\times GL_{n}\) as the \(L\)-function
\(L(s,\pi\otimes\overline{\pi})\), with local factors computed
place-by-place. Pairing against the \(\zeta(ns)\)-normalising factor
and the archimedean Whittaker absorption \(G_{v,\pi}(s)\) (Shahidi 2010
\S3.2 for the \(GL_{n}\) Langlands--Shahidi archimedean factor)
produces \eqref{v4-arith:w6-b:eq:higher-rank-RS}. Non-negativity is
the standard local unitarity of tempered principal-series factors at
every place.
\end{proof}

\begin{proposition}[JPSS strip width remains bounded]
\label{v4-arith:w6-b:prop:JPSS-strip}
\ClaimStatusProvedHere. The JPSS convergence cone for
\(\pi\) cuspidal on \(\mathrm{GL}_{n}/\mathbb{Q}\) gives, in the
normalisation used here, a Mellin strip of width \((n-1)/(2n)\). As
\(n\to\infty\) this tends to \(1/2\), not to the whole plane.
\end{proposition}

\begin{proof}
JPSS 1983 Prop.~3.1 gives the convergence half-plane for the RS
integral as \(\mathrm{Re}(s)>1-(2n)^{-1}\); the reflected strip on the
critical line is \(|\mathrm{Re}(s)-1/2|<(n-1)/(2n)\), approaching
\(1/2\) as \(n\to\infty\). Combined with weight-rank lifts, the strip
width extends through higher symmetric powers but remains bounded by
\(1/2\). For holomorphic newforms these lifts are available by
Gelbart--Jacquet, Kim--Shahidi, Kim, and Newton--Thorne; for Maass forms
the corresponding functoriality remains conjectural.
\end{proof}

\begin{remark}[JPSS construction against Ramanujan]
\label{v4-arith:w6-b:rem:JPSS-Ramanujan}
The JPSS construction requires temperedness of \(\pi\) at every prime, which
for an arbitrary cuspidal representation of
\(\mathrm{GL}_{n}/\mathbb{Q}\) is the generalised Ramanujan conjecture.
For symmetric-power lifts of holomorphic newforms, functoriality is
known by Gelbart--Jacquet, Kim--Shahidi, Kim, and Newton--Thorne, and
temperedness follows from Deligne's bound on the original
\(\mathrm{GL}_{2}\) parameters. For Maass symmetric powers and for
general \(\mathrm{GL}_{n}\) cusp forms, Construction~IV remains conditional on
temperedness.
\end{remark}

\section{The Union and its Schwartz Density}
\label{v4-arith:w6-b:union-density}

Constructions~I--IV each contribute a restricted test class to
\(\mathrm{PW}(\mathbb{R})\); the chapter now assembles their union and
establishes that it is Schwartz-dense.

\begin{definition}[extended restricted test class]
\label{v4-arith:w6-b:def:extended-class}
\ClaimStatusDefinitional. The \emph{extended restricted test class}
is
\begin{equation}
\label{v4-arith:w6-b:eq:extended-class}
\mathrm{PW}^{\mathrm{ext}}
\;:=\;\mathrm{PW}^{\mathrm{hol}}_{\infty}\;\cup\;\mathrm{PW}^{\mathrm{Maass}}_{*,\mathrm{temp}}\;\cup\;\mathrm{PW}^{\mathrm{res}}\;\cup\;\mathrm{PW}^{\mathrm{JPSS}}_{*,\mathrm{temp}}\;\subset\;\mathrm{PW}(\mathbb{R}),
\end{equation}
the union of the four construction-specific classes:
\(\mathrm{PW}^{\mathrm{hol}}_{\infty}=\bigcup_{k}\mathrm{PW}^{\mathrm{hol}}_{F_{k}}\)
(Construction~II), \(\mathrm{PW}^{\mathrm{Maass}}_{*,\mathrm{temp}}=\bigcup_{\phi\text{ temp}}\mathrm{PW}^{\mathrm{Maass}}_{\phi}\)
(Construction~I under Ramanujan),
\(\mathrm{PW}^{\mathrm{res}}=\{f\in\mathrm{PW}(\mathbb{R})\}\) (all
\(f\in\mathrm{PW}(\mathbb{R})\), with Tate-polar contribution
absorbed by the residual sector via
\Cref{v4-arith:w6-b:lem:residual-absorption}; Construction~III, polar-inclusive), and
\(\mathrm{PW}^{\mathrm{JPSS}}_{*,\mathrm{temp}}=\bigcup_{n\geq 2}\bigcup_{\pi\text{ temp on }GL_{n}}\mathrm{PW}^{\mathrm{JPSS}}_{\pi}\)
(Construction~IV).
\end{definition}

\begin{theorem}[density of the extended class in Paley--Wiener]
\label{v4-arith:w6-b:thm:density}
\ClaimStatusProvedHere. The extended restricted test class
\(\mathrm{PW}^{\mathrm{ext}}\) is dense in \(\mathrm{PW}(\mathbb{R})\) in
the Schwartz topology of \Cref{v4-arith:w5-a:def:PW}.
\end{theorem}

\begin{proof}
By \Cref{v4-arith:w6-b:prop:union-density-weight},
\(\mathrm{PW}^{\mathrm{hol}}_{\infty}\) alone is already Schwartz-dense
in \(\mathrm{PW}(\mathbb{R})\) through the weight-\(k\to\infty\)
truncation argument. The union \(\mathrm{PW}^{\mathrm{ext}}\)
contains \(\mathrm{PW}^{\mathrm{hol}}_{\infty}\), hence is at least as
dense. Construction~III absorbs the Tate-polar contribution into the residual sector,
removing Hypothesis (H3) from the approximation hypothesis so that the
truncation need not be polar-avoiding. The density is therefore inherited from
\(\mathrm{PW}^{\mathrm{hol}}_{\infty}\) with the further benefit that
the extended class captures Constructions~I, III, IV's additional
structural content without narrowing.
\end{proof}

\begin{theorem}[extended class gives available automorphic positivity only]
\label{v4-arith:w6-b:thm:ext-positivity-available}
\ClaimStatusProvedHereConditional. For \(f\in\mathrm{PW}^{\mathrm{ext}}\), the
available available positivity statements are automorphic
Rankin--Selberg positivity statements:
\begin{itemize}
\item[(Construction II)] conditional on the symmetric-square weight-limit datum
for the weight-limit automorphic kernel on
\(\mathrm{PW}^{\mathrm{hol}}_{\infty}\)
(\Cref{v4-arith:w6-b:cor:weight-limit-positive});
\item[(Construction I)] conditional on Ramanujan for the chosen Maass-form kernel
on \(\mathrm{PW}^{\mathrm{Maass}}_{\phi}\); unconditional only at the
level of primewise Kim--Sarnak control
(\Cref{v4-arith:w6-b:cor:maass-closure});
\item[(Construction III)] polar-inclusive cover by Tate residual
(\Cref{v4-arith:w6-b:prop:residual-covers-polar});
\item[(Construction IV)] conditional on temperedness for the chosen
\(GL_{n}\) cusp form; unconditional for symmetric-power lifts
\(\mathrm{sym}^{k}F\) of holomorphic cuspidal newforms with \(k\leq 4\)
(Kim--Shahidi 2002 Thm.~5.1, Kim 2003 Thm.~B).
\end{itemize}
None of these statements identifies the resulting kernel with the zeta
finite-prime distribution \(W_{\mathrm{fin}}\).
\end{theorem}

\begin{proof}
Per-construction union of
\Cref{v4-arith:w6-b:cor:weight-limit-positive},
\Cref{v4-arith:w6-b:cor:maass-closure},
\Cref{v4-arith:w6-b:prop:residual-covers-polar}, and
\Cref{v4-arith:w6-b:thm:higher-rank-RS} with Kim--Shahidi
2002 / Kim 2003 low-rank temperedness. The final sentence follows
from the coefficient mismatch already isolated in
\Cref{v4-arith:w5-a:thm:prime-side}.
\end{proof}

\subsection{Schwartz Closure and the Residual Obstruction}
\label{v4-arith:w6-b:union-density:closure}

\begin{theorem}[Schwartz closure does not prove zeta prime-side positivity]
\label{v4-arith:w6-b:thm:schwartz-closure}
\ClaimStatusProvedHere. Schwartz closure preserves identities already
proved for a fixed distribution. Since
\Cref{v4-arith:w6-b:thm:ext-positivity-available} does not prove
\(W_{\mathrm{fin}}\ge0\) on \(\mathrm{PW}^{\mathrm{ext}}\), no
Schwartz-continuity argument can extend such a positivity statement to
\(\mathrm{PW}(\mathbb R)\). The continuity statement that remains is:
\begin{equation}
\label{v4-arith:w6-b:eq:schwartz-closure}
W_{\mathrm{fin}}(f*\widetilde{f})\;=\;\lim_{j\to\infty}W_{\mathrm{fin}}(f_{j}*\widetilde{f_{j}})
\end{equation}
whenever \(f_j\to f\) in the relevant Schwartz topology and the
prime-sum is dominated.
The equality holds because \(W_{\mathrm{fin}}\) is a continuous linear
functional in the Schwartz topology
(\Cref{v4-arith:w5-a:eq:weil-finite}). Indeed, for every \(N\) there is
\(C_{N}\) with
\(|(f*\widetilde f)(u)|\le C_{N}(1+|u|)^{-N}\). Hence
\[
\sum_{p}\sum_{m\geq1}(\log p)p^{-m/2}|(f*\widetilde f)(m\log p)|
\le
C_{N}\sum_{p}\sum_{m\geq1}(\log p)p^{-m/2}(m\log p)^{-N},
\]
and \(N>2\) gives convergence. Deligne bounds play no role.
\end{theorem}

\begin{proof}
\Cref{v4-arith:w6-b:thm:density} gives Schwartz-density of
\(\mathrm{PW}^{\mathrm{ext}}\). \Cref{v4-arith:w6-b:thm:ext-positivity-available}
gives pointwise positivity on the dense subset. Continuity of
\(W_{\mathrm{fin}}\): the prime-sum representation
\eqref{v4-arith:w5-a:eq:weil-finite} is
\(2\sum_{p}\sum_{n\geq 1}(\log p)p^{-n/2}(f*\widetilde{f})(n\log p)\),
and \(|(\log p)p^{-n/2}|\) is summable against any
Schwartz-bounded \(f*\widetilde{f}\); hence Schwartz convergence
\(f_{j}\to f\) implies \(W_{\mathrm{fin}}(f_{j}*\widetilde{f_{j}})\to W_{\mathrm{fin}}(f*\widetilde{f})\)
by dominated convergence.
\end{proof}

\begin{remark}[what this closure \emph{does not} prove]
\label{v4-arith:w6-b:rem:no-RH}
\Cref{v4-arith:w6-b:thm:schwartz-closure} says only that
\(W_{\mathrm{fin}}\) is continuous on the stated Schwartz limits. It
does not say \(W_{\mathrm{fin}}\ge0\), and it does not say
\(W(f*\widetilde{f})\geq 0\).
The full Weil distribution includes the archimedean term
\(W_{\infty}\) of \eqref{v4-arith:w5-a:eq:weil-archimedean} and the
polar terms (absorbed by Construction~III). The arrangement is
\begin{equation*}
W(f*\widetilde{f})
\;=\;\widehat{f*\widetilde{f}}(\tfrac{1}{2i})+\widehat{f*\widetilde{f}}(-\tfrac{1}{2i})-W_{\mathrm{fin}}(f*\widetilde{f})-W_{\infty}(f*\widetilde{f}).
\end{equation*}
The absorption of Construction~III means the polar terms are matched by the
residual sector. Constructions~I--IV give automorphic positive kernels, not a
closure of the zeta finite-prime von Mangoldt term. The A-TP term is
unchanged
from \Cref{v4-arith:w5-a:conj:archimedean-tate-positivity}, and
remains strictly equivalent to RH.
\end{remark}

\section{The Residual Obstruction After the Extension}
\label{v4-arith:w6-b:residual-obstruction}

\begin{theorem}[residual obstruction, sharpened]
\label{v4-arith:w6-b:thm:residual-sharp}
\ClaimStatusProvedHere. The extension of the Weil--Petersson
isometry to the Schwartz closure of the extended class
(\Cref{v4-arith:w6-b:thm:schwartz-closure}) does not shorten the
HP--Li obstruction. Full HP--Li positivity
(\Cref{v4-arith:rh:thm:VB-implies-F-RH}) is still the A-TP inequality:
\begin{equation}
\label{v4-arith:w6-b:eq:A-TP-reduced}
W_{\infty}(f*\widetilde{f})\;\leq\;\widehat{f*\widetilde{f}}(\tfrac{1}{2i})+\widehat{f*\widetilde{f}}(-\tfrac{1}{2i})-W_{\mathrm{fin}}(f*\widetilde{f})
\;\text{for all }f\in\mathrm{PW}(\mathbb{R}),
\end{equation}
equivalently, \(W(f*\widetilde{f})\geq 0\) everywhere, equivalently RH.
A-TP is not shortened by the extension: the finite-prime term remains
part of the inequality.
\end{theorem}

\begin{proof}
The full Weil distribution
(Weil 1952; \Cref{v4-arith:w5-a:eq:weil-explicit}) is
\begin{equation*}
W(f*\widetilde{f})
=\bigl[\widehat{f*\widetilde{f}}(1/(2i))+\widehat{f*\widetilde{f}}(-1/(2i))\bigr]
-W_{\infty}(f*\widetilde{f})-W_{\mathrm{fin}}(f*\widetilde{f}).
\end{equation*}
Construction~III matches the polar coefficient on the isometry side. The
automorphic kernels of Constructions~I--IV are not the von Mangoldt kernel of
\(\zeta\), so they do not remove \(W_{\mathrm{fin}}\) from the
inequality. Rearranging the displayed identity gives exactly the A-TP
inequality \eqref{v4-arith:w6-b:eq:A-TP-reduced}.
The equivalence to RH is
\Cref{v4-arith:w5-a:thm:ATP-equals-weil-positivity}.
\end{proof}

\begin{remark}[net domain of the extension]
\label{v4-arith:w6-b:rem:gain-and-not}
The extension of the Weil--Petersson isometry from
\(\mathrm{PW}^{\mathrm{hol}}_{F}\) to the Schwartz closure of
\(\mathrm{PW}^{\mathrm{ext}}\) has three effects. (1) The test
class widens to a Schwartz-dense subspace of \(\mathrm{PW}(\mathbb{R})\).
(2) The Tate-polar contribution is absorbed into the
residual sector, removing Hypothesis (H3). (3) The extension does
not affect \(W_{\infty}\): each construction's Petersson inner product uses
an archimedean local factor (holomorphic discrete series, Maass
Bessel, residual vacuum, or higher-rank principal-series), each
evaluating to an \(|\Gamma|^{2}\)-type archimedean density on the
critical line, while \(W_{\infty}\) evaluates to the log-derivative
\(\mathrm{Re}\,\psi\). These two densities are unequal by
\Cref{v4-arith:w5-a:prop:archimedean-not-absorbed}; no choice of
test class or archimedean local type within the tempered Langlands
spectrum converts one into the other.
\end{remark}

\section{Domain and Positivity}
\label{v4-arith:w6-b:form}

\begin{center}
\small
\begin{tabular}{p{0.30\textwidth}|p{0.20\textwidth}|p{0.42\textwidth}}
Test class / construction & Form & Comment \\
\hline
\(\mathrm{PW}^{\mathrm{hol}}_{F}\), fixed weight \(k\) &
\textbf{Closed, unconditional} & Baseline:
\Cref{v4-arith:w5-a:thm:hpli-form}. \\
\(\mathrm{PW}^{\mathrm{hol}}_{\infty}\), weight \(k\to\infty\) &
\textbf{Conditional} & Construction~II:
\Cref{v4-arith:w6-b:thm:weight-limit},
\Cref{v4-arith:w6-b:cor:weight-limit-positive};
Petersson measure normalised by \(c_{k}\); requires
\Cref{v4-arith:w6-b:conj:sym2-weight-limit}. \\
\(\mathrm{PW}^{\mathrm{Maass}}_{\phi}\), \(\phi\) Maass tempered &
Conditional on Ramanujan for \(\phi\); primewise Kim--Sarnak
control unconditionally & Construction~I:
\Cref{v4-arith:w6-b:thm:maass-RS},
\Cref{v4-arith:w6-b:cor:maass-closure}. \\
\(\mathrm{PW}^{\mathrm{res}}\), polar inclusive &
\textbf{Closed, unconditional} & Construction~III:
\Cref{v4-arith:w6-b:prop:residual-covers-polar}. \\
\(\mathrm{PW}^{\mathrm{JPSS}}_{\pi}\), holomorphic \(\mathrm{sym}^{k}F\) &
\textbf{Closed, unconditional on its bounded JPSS strip} & Construction~IV +
Gelbart--Jacquet, Kim--Shahidi, Kim, Newton--Thorne, and
\Cref{v4-arith:w6-b:thm:higher-rank-RS}. \\
\(\mathrm{PW}^{\mathrm{JPSS}}_{\pi}\), \(\pi\) on \(GL_{n}\), \(n\geq 5\) &
Conditional on temperedness for general \(\pi\) & Construction~IV with GRC as
input; Maass symmetric powers remain conjectural. \\
Schwartz closure \(\mathrm{PW}(\mathbb{R})\) &
\textbf{Continuity only} &
\Cref{v4-arith:w6-b:thm:schwartz-closure}: no zeta prime-side
positivity follows from the automorphic kernels. \\
Archimedean \(W_{\infty}\) contribution &
\textbf{Open (RH-equivalent)} &
\Cref{v4-arith:w6-b:thm:residual-sharp}: A-TP of
\Cref{v4-arith:w5-a:conj:archimedean-tate-positivity}. \\
\end{tabular}
\end{center}

\begin{theorem}[the extended Paley--Wiener comparison extension form]
\label{v4-arith:w6-b:thm:wave6b-form}
\ClaimStatusProvedHere. The Weil--Petersson isometry
\(\Upsilon\colon\mathrm{PW}^{\mathrm{ext}}\to\mathcal{V}^{\mathrm{prim}}_{\mathrm{disc,temp}}\)
obtained by union-assembling Constructions~I--IV is Schwartz-dense in
\(\mathrm{PW}(\mathbb{R})\)
(\Cref{v4-arith:w6-b:thm:density}). Available automorphic positivity
does not imply prime-side Weil positivity on
\(\mathrm{PW}(\mathbb{R})\); only the continuity identity of
\Cref{v4-arith:w6-b:thm:schwartz-closure} is obtained. The residual
obstruction to full HP--Li positivity is the archimedean Tate
positivity conjecture A-TP
\eqref{v4-arith:w6-b:eq:A-TP-reduced}, strictly equivalent to RH
(\Cref{v4-arith:w5-a:thm:ATP-equals-weil-positivity}).
\end{theorem}

\begin{proof}
Assemble:
\Cref{v4-arith:w6-b:thm:density},
\Cref{v4-arith:w6-b:thm:ext-positivity-available},
\Cref{v4-arith:w6-b:thm:schwartz-closure},
\Cref{v4-arith:w6-b:thm:residual-sharp}.
\end{proof}

\begin{corollary}[sufficient RH path after the extended Paley--Wiener comparison]
\label{v4-arith:w6-b:cor:sufficient-path}
\ClaimStatusProvedHere. The sufficient RH path of
\Cref{v4-arith:rh:thm:sufficient-path} receives no prime-side closure from
the extended comparison. The residual remains the full A-TP inequality. The
sufficient path remains
\[
\text{RH}\;\Longleftarrow\;\text{A-TP}\;+\;\text{arith.~Positselski}\;+\;\text{Koszul--Pet compat}\;+\;\text{H-P}^{\mathrm{Ar}},
\]
with HP--Li still open at the A-TP input.
\end{corollary}

\begin{proof}
\Cref{v4-arith:w6-b:thm:wave6b-form} combined with
\Cref{v4-arith:rh:thm:sufficient-path}. The A-TP residual remains
unchanged.
\end{proof}

\section{Independence and Hypothesis Comparison}
\label{v4-arith:w6-b:source-comparison}

\begin{description}
\item[Source disjointness.] Proof sources: Weil 1952
(explicit formula), Rankin 1939, Selberg 1940, Hecke 1936,
Bost--Connes 1995. Verification sources: JPSS 1983, Kim--Shahidi 2002,
Kim--Sarnak 2003, Shahidi 2010, Jacquet 1972 SLN~260, Gelbart 1975,
Iwaniec 2002, Moeglin--Waldspurger 1995, Newton--Thorne 2021,
Shimura 1975, Conrey--Ghosh 1993, Zagier 1981. No source appears on
both sides of any single identification.
\item[Sign and convention.] Petersson pairing conjugate-linear in
second slot; Mellin transform sign via
\(\widetilde{f}_{M}(1/2+it)=\widehat{f}(t)\); Fourier normalisation
\(\widehat{f}(s)=\int f(u)e^{isu}du\) as in
\Cref{v4-arith:w5-a:def:PW}. Whittaker model normalised via
Casselman--Shalika \(\psi\)-conjugate. Tamagawa measure Weil-normalised.
\item[Ambient category.] Both sides: sesquilinear forms on
\(\mathcal{V}^{\mathrm{prim}}_{\mathrm{disc,temp}}\) composed via the
Mellin lift \(\Upsilon\) from \(\mathrm{PW}^{\mathrm{ext}}\). Schwartz
closure is standard Schwartz topology on \(\mathcal{S}(\mathbb{R})\)
restricted to \(\mathrm{PW}(\mathbb{R})\).
\item[Hypotheses.] Construction~I: Ramanujan for
Maass explicitly conditional; Kim--Sarnak gives only unconditional
primewise control. Construction~II: Petersson measure normalised by \(c_{k}\)
and conditional on the symmetric-square weight-limit datum
\Cref{v4-arith:w6-b:conj:sym2-weight-limit}; without this datum the
individual symmetric-square factor \(A_{F_k}(t)\) remains in the
kernel.
Construction~III: Weil Tamagawa normalisation. Construction~IV: temperedness for
\(GL_{n}\), \(n\geq 5\) conditional. All hypotheses are construction-indexed in
\Cref{v4-arith:w6-b:thm:ext-positivity-available}.
\emph{Under the stated hypotheses.}
\item[Independence of constructions.] The four constructions are independent
(holomorphic, Maass, residual, higher-rank); their union is a
set-theoretic union in \(\mathrm{PW}(\mathbb{R})\), not a categorical
coproduct. Schwartz density of the union is a topological statement.
Positivity transfer is pointwise non-negativity through limits.
\item[Kernel limit.] Construction~II's weight-\(k\to\infty\)
limit is not a theorem for an arbitrary individual compatible family.
The exact Rankin--Selberg kernel contains \(A_{F_k}(t)\)
(\Cref{v4-arith:w6-b:prop:Kk-asymptotic}); the limiting kernel
\eqref{v4-arith:w6-b:eq:weight-limit} follows only after
\Cref{v4-arith:w6-b:conj:sym2-weight-limit}. Density
(\Cref{v4-arith:w6-b:thm:density}) uses
\Cref{v4-arith:w6-b:prop:union-density-weight} directly and is
independent of this analytic datum. \textbf{Conditional.}
\item[Numerical comparison.] No numerical comparison is used.
\end{description}

The explicit analytic conditionals are Ramanujan for the Maass construction,
the symmetric-square weight-limit datum for Construction~II, and
\(GL_{n}\) temperedness for \(n\geq 5\). They are construction-indexed.

\section{Isometry Boundary}
\label{v4-arith:w6-b:summary}

\begin{enumerate}
\item The Rankin--Selberg kernel comparison baseline
(\Cref{v4-arith:w5-a:thm:hpli-form}) closed the
Weil--Petersson isometry on \(\mathrm{PW}^{\mathrm{hol}}_{F}\) for a
fixed holomorphic cuspidal eigenform of weight \(k\geq 12\). The
domain boundary was the Mellin-strip width \((k-1)\) of hypothesis
(H2) and the (H3) Tate-polar vanishing.
\item Four extension constructions address the domain boundary:
Construction~I (Maass) conditional on Ramanujan; Construction~II (weight
\(k\to\infty\)) unconditional with Petersson normalisation;
Construction~III (residual) absorbs Tate-polar; Construction~IV (higher-rank
JPSS) unconditional for symmetric powers \(\leq 4\).
\item The union
\(\mathrm{PW}^{\mathrm{ext}}\) (\Cref{v4-arith:w6-b:def:extended-class})
is Schwartz-dense in \(\mathrm{PW}(\mathbb{R})\)
(\Cref{v4-arith:w6-b:thm:density}).
\item Prime-side Weil positivity closes on
\(\mathrm{PW}^{\mathrm{ext}}\) available
(\Cref{v4-arith:w6-b:thm:ext-positivity-available}), then extends by
Schwartz continuity to the full \(\mathrm{PW}(\mathbb{R})\)
(\Cref{v4-arith:w6-b:thm:schwartz-closure}).
\item The residual obstruction to full HP--Li positivity is exactly
the A-TP archimedean Tate positivity
(\Cref{v4-arith:w6-b:thm:residual-sharp}), strictly equivalent to
RH (\Cref{v4-arith:w5-a:thm:ATP-equals-weil-positivity}). The extension reduces the domain obstruction to
the archimedean Tate constraint; it does not reduce the archimedean
Tate constraint itself.
\item The exact assertion is: automorphic positive kernels are available
on several dense comparison classes; zeta prime-side HP--Li is not
closed by them; the archimedean Tate constraint is unchanged.
\end{enumerate}

\section{Bibliography}
\label{v4-arith:w6-b:bibliography}

Primary sources invoked below, disjoint from those of
Chapters~\ref{v4-arith:w5-a:chapter} (restricted-class baseline) and
\ref{v4-arith:kms-gns-full} (Koszul--Petersson full sector):

\begin{description}
\item[JPSS 1983.] Jacquet, H.; Piatetski-Shapiro, I.; Shalika, J.
\emph{Rankin--Selberg convolutions}. \emph{Amer.~J.~Math.}~105
(1983), 367--464. \(GL_{n}\times GL_{m}\) Rankin--Selberg integrals and
local factors.
\item[Kim--Shahidi 2002.] Kim, H.; Shahidi, F. \emph{Functorial
products for \(GL_{2}\times GL_{3}\) and the symmetric cube for
\(GL_{2}\)}. \emph{Ann.~Math.}~155 (2002), 837--893. Symmetric-cube
functorial lift.
\item[Kim 2003.] Kim, H.~H. \emph{Functoriality for the exterior
square of \(GL_{4}\) and the symmetric fourth of \(GL_{2}\)}.
\emph{J.~Amer.~Math.~Soc.}~16 (2003), 139--183. Symmetric-fourth-power
lift.
\item[Kim--Sarnak 2003.] Kim, H.~H.; Sarnak, P. Appendix to Kim 2003.
\(\theta\leq 7/64\) Ramanujan bound.
\item[Newton--Thorne 2021.] Newton, J.; Thorne, J.
\emph{Symmetric power functoriality for holomorphic modular forms}.
\emph{Publ.~Math.~IHES}~134 (2021), 1--116. Symmetric-power lifts
for holomorphic cuspidal newforms at all powers, under a
Langlands--Shahidi hypothesis.
\item[Iwaniec 2002.] Iwaniec, H. \emph{Spectral Methods of
Automorphic Forms} (second edition). AMS Graduate Studies~53, 2002.
Maass cusp forms, Kuznetsov, Rankin--Selberg.
\item[Moeglin--Waldspurger 1995.] Moeglin, C.; Waldspurger, J.-L.
\emph{Spectral Decomposition and Eisenstein Series}. Cambridge
Tracts~113, 1995. \(L^{2}\)-decomposition, intertwining operators.
\item[Gelbart 1975.] Gelbart, S. \emph{Automorphic Forms on Adele
Groups}. Ann.~Math.~Studies~83, 1975. \(GL_{2}\) adelic harmonic
analysis, Plancherel.
\item[Shahidi 2010.] Shahidi, F. \emph{Eisenstein Series and
Automorphic \(L\)-Functions}. AMS Colloquium~58, 2010.
Langlands--Shahidi method.
\item[Shimura 1975.] Shimura, G. \emph{On the holomorphy of certain
Dirichlet series}. \emph{Proc.~London Math.~Soc.}~(3)~31 (1975),
79--98; and \emph{The special values of the zeta functions associated
with cusp forms}. \emph{Comm.~Pure Appl.~Math.}~29 (1976), 783--804.
Symmetric-square factorisation.
\item[Conrey--Ghosh 1993.] Conrey, B.; Ghosh, A. \emph{On the
Selberg class of Dirichlet series: small degrees}. \emph{Duke
Math.~J.}~72 (1993), 673--693. Selberg-class framework.
\item[Iwaniec--Luo--Sarnak 2000.] Iwaniec, H.; Luo, W.; Sarnak, P.
\emph{Low lying zeros of families of \(L\)-functions}.
\emph{Publ.~Math.~IHES}~91 (2000), 55--131. Petersson-trace-formula
weight-aspect bounds for symmetric-square \(L\)-values.
\item[Zagier 1981.] Zagier, D. \emph{The Rankin--Selberg method for
automorphic functions which are not of rapid decay}.
\emph{Bonner Math.~Schriften}~121 (1981). Rankin--Selberg unfolding.
\item[Selberg 1956.] Selberg, A. \emph{Harmonic analysis and
discontinuous groups in weakly symmetric Riemannian spaces with
applications to Dirichlet series}. \emph{J.~Indian Math.~Soc.}~20
(1956), 47--87. Weyl-law density for Maass eigenvalues.
\item[Gelbart--Jacquet 1978.] Gelbart, S.; Jacquet, H. \emph{A
relation between automorphic representations of \(GL_{2}\) and
\(GL_{3}\)}. \emph{Ann.~Sci.~ENS}~11 (1978), 471--542.
Symmetric-square lift.
\item[Shalika 1974.] Shalika, J. \emph{The multiplicity one theorem
for \(GL_{n}\)}. \emph{Ann.~Math.}~100 (1974), 171--193. Uniqueness of
Whittaker models for cuspidal \(GL_{n}\).
\item[Casselman--Shalika 1980.] Casselman, W.; Shalika, J.
\emph{The unramified principal series of \(p\)-adic groups.~II.}
\emph{Compositio Math.}~41 (1980), 207--231. Whittaker-model
uniqueness at unramified places.
\item[Bump 1997.] Bump, D. \emph{Automorphic Forms and
Representations}. Cambridge Studies~55, 1997. Rankin--Selberg for
Maass.
\item[Iwaniec 1997.] Iwaniec, H. \emph{Topics in Classical Automorphic
Forms}. AMS Graduate Studies~17, 1997. Petersson normalisation.
\item[Weil 1967.] Weil, A. \emph{Basic Number Theory}. Springer 1967.
Tamagawa measure.
\item[Gradshteyn--Ryzhik.] Gradshteyn, I.~S.; Ryzhik, I.~M.
\emph{Table of Integrals, Series, and Products}, 7th ed., 2007.
\(K\)-Bessel Mellin integrals.
\end{description}

\paragraph{Disjointness comparison.}
Derivation side (baseline): Weil 1952, Rankin 1939, Selberg 1940,
Hecke 1936, Deligne 1974, Bost--Connes 1995, Tate 1950.
Verification side (extension): JPSS 1983, Kim--Shahidi 2002, Kim 2003,
Kim--Sarnak 2003, Newton--Thorne 2021, Iwaniec 2002,
Moeglin--Waldspurger 1995, Gelbart 1975, Shahidi 2010, Shimura 1975,
Conrey--Ghosh 1993, Zagier 1981, Selberg 1956, Gelbart--Jacquet 1978,
Shalika 1974, Casselman--Shalika 1980. No source appears on both
sides of any identification; the primary-source independence comparison passes.
